% ==================== 第 6 章 总体结果、失败分析与适用边界 ====================
\section{总体结果、失败分析与适用边界}

\subsection{本作品实际总体表现：代表查询与内置样例库成绩}

\noindent\textbf{两个代表查询的完整档案}（口径与第 1 章核心主张一致）：

\begin{longtable}{@{}p{3.2cm}p{2.6cm}p{3.6cm}p{2.6cm}p{2.4cm}@{}}
\toprule
实际任务或来源范围 & 实际处理的记录、条目或数据量 & 主要质量结果 & 能否找到原始出处并继续使用 & 仍存在的问题 \\
\midrule
疏散星团 M45（年龄、距离、距离模数、视差、金属丰度；29 篇论文 + MWSC 星团目录，1998--2025，10 种期刊） & 103 条记录 / 78 张图证 / 69 条修改留痕 & 综合质量分 0.85；现代测距值收敛 134.4--136.2 pc、年龄 100--130 Myr、金属丰度 $-$0.03~+0.03 dex；31 项跨源冲突全部保留并归因 & 是（99 条 96.1\% 绑定页码与定位框，可逐条回溯原文） & 范围写法记录（如 “75-100 Myr”）按设计保留待核对；跨来源单位差异无法保证全局统一 \\
单一恒星 Pollux（质量、半径、表面温度、光谱型；7 篇文献 + 4 个星表） & 16 条记录 / 6 处修改留痕 & 综合质量分 0.9235（Excellent）；零未解冲突、免人工复核 & 是（记录全部带来源与原文位置） & 星表缺失字段（质量、半径）保留显式缺失标记，未填充 \\
\bottomrule
\end{longtable}

\noindent\textbf{内置样例库：17 组真实查询的运行成绩单}（与前端演示库同源；质量分为系统自动评估产出，取自任务状态与结果端点同源数据；覆盖恒星个体、双星与致密天体、星团与星协、星系与活动星系核、星系团等对象类别与约 40 种物理性质）：

\begin{longtable}{@{}p{3.2cm}p{4.8cm}p{1.7cm}p{1.7cm}p{2.0cm}@{}}
\toprule
样例对象 & 查询性质 & 记录数 & 来源数 & 自动质量分 \\
\midrule
参宿四（红超巨星） & 半径、质量、表面温度 & 47 & 34 & 0.9024 \\
大角星（红巨星） & 距离、金属丰度、有效温度 & 52 & 18 & 0.8144 \\
Pollux（巨星，点源） & 质量、半径、表面温度、光谱型 & 16 & 11 & 0.9235 \\
天狼星 B（白矮星） & 质量、表面重力、冷却年龄 & 31 & 11 & 0.8908 \\
天鹅座 X-1（黑洞双星） & 轨道周期、致密天体质量 & 48 & 28 & 0.9153 \\
$\delta$ Cephei（造父变星） & 光变周期、距离、金属丰度 & 53 & 31 & 0.7327 \\
比邻星（红矮星） & 距离、视差、自行 & 50 & 16 & 0.7751 \\
昴星团 M45（08-28 定稿版） & 年龄、距离、金属丰度（→5 性质） & 103 & 30 & 0.8502 \\
昴星团 M45（09-01 复跑版） & 同上 & 111 & 28 & 0.8651 \\
半人马座 $\omega$（球状星团） & 中心速度弥散、总质量、金属丰度 & 25 & 24 & 0.9085 \\
五车二（双星） & 轨道周期、双星质量、视向速度半振幅 & 23 & 8 & 0.9142 \\
M87（活动星系核） & 黑洞质量、爱丁顿比、6cm 射电流量 & 34 & 22 & 0.9345 \\
NGC 4151（赛弗特星系） & 黑洞质量、H$\beta$ 线半高全宽、X 射线谱指数 & 99 & 28 & 0.8507 \\
3C 273（耀变体） & 红移、绝对星等 & 15 & 13 & 0.9013 \\
HD 209458b（系外行星） & 轨道周期、行星质量、行星半径 & 45 & 36 & 0.8766 \\
后发座星系团 & 速度弥散、X 射线温度、M500 & 19 & 11 & 0.9212 \\
SN 2011fe（Ia 型超新星） & 峰值绝对星等 & 32 & 10 & 0.9550 \\
\bottomrule
\end{longtable}

\noindent 说明：① 库内 M45 有两个真实运行版本，正文旗舰口径统一取 08-28 定稿版（103 条 / 0.85 / 30 源，即 1.2 结果 1）；② 质量分分布如实展示（0.73--0.96），不全部是高分的“精选观感”——演示库的选取标准是“代表性+性质齐全”，而非分数筛选；③ 每条记录均可点击回溯原文（溯源能力见 3.5）。

\subsection{真实失败与边界}

\begin{longtable}{@{}p{2.6cm}p{4.6cm}p{3.8cm}p{3.4cm}@{}}
\toprule
典型失败或边界问题 & 实际表现 & 原因分析 & 当前处理方式或支持程度 \\
\midrule
\textbf{多星系统的组分混淆（天狼星 A/B）} & 真实查询“天狼星 B 的质量、有效温度、表面重力和冷却年龄”（任务 d15a54ae，36 条，自动质量分 0.89）：effective\_temperature 的 8 条中 4 条是 \textbf{A 星值}（Gaia DR2 的 9267/7245/9726 K）与 1 条错误值（51111 K），仅 3 条正确（约 25000 K）；质量另有 2 条混入前身星与 IFMR 他星值 & 多星系统论文与星表条目中 A/B 子星数值并存，提取阶段缺乏“数值归属哪颗子星”的成员级确认；单值均在物理合理范围，自动质检无法拦截 & 系统侧尚不能自动判定子星归属（属开放问题，人工复核可发现但无法全自动裁决）；团队以查询设计规避——演示库改用不含该性质的重查版（31 条，3/3 干净），历史任务原始记录保留并标注，不静默删除 \\
\textbf{原型类名天体的种群污染（大陵五 / 米拉）} & 同批次真实查询：大陵五（45 条）orbital\_period 仅 4/14 为本星（2.8673 d），10 条混入 SZ Herculis、IZ Tel/UW Vir 等其他双星周期，component\_mass 20 条中 17 条为他星或候选伴星值；米拉（41 条）pulsation\_period 仅 4/23 为本星，19 条为同文\textbf{其他 Mira 星}（2021OEJV 种群统计论文） & 以“类名原型天体”（Algol 型 / Mira 型）检索命中的文献主体是“该类型恒星”的统计研究，模型把同文其他成员星的数值提取为目标值——\textbf{值正确但对象错误}，且单值均在合理范围，自动质检无法区分 & 这是当前最系统的失败模式（与历史 RR Lyr 同类）；目前以“论文聚焦目标天体本身”的样例筛选缓解，系统侧需要论文级目标确认能力（判断文献研究对象是否=查询对象），尚不支持 \\
\bottomrule
\end{longtable}

\noindent 上述两例与 1.2 局限三条（版权/性质库覆盖/评估客观性）、4.3 未解决问题（范围写法保留、单位差异无法全局统一）共同构成系统的诚实边界；逐条核对存档见 \code{docs/demo-query-analysis}，随附录提交。

\subsection{应用与成本综合说明}

\begin{itemize}[leftmargin=1.2em,itemsep=3pt]
\item \textbf{结构化数据已如何被下载、分析程序或后续科研流程使用：}每次查询产出的长表/宽表 CSV 与结构化 JSON 可直接导入 Python（pandas）、R 等分析环境——研究者无需再与“文本混排的论文数值”搏斗，即可完成参数汇总、误差传播与多文献对比；带置信度与质量分的记录支持按阈值筛选子集（如取高置信记录做恒星演化模型标定或距离尺度拟合），冲突标注与测量方法列则使“文献值汇编数据集”（附方法学差异）可被综述写作与后续检验直接引用；溯源元数据（DOI/BibTeX、页码与定位框）可一键对接文献管理与论文写作流程，无需二次核对出处；多次查询的结果还可跨天体拼装成样本级数据集，供统计研究复用。系统输出的“值+单位+出处+状态”四要素齐备，是进入下游科研管线的可直接消费形态。
\item \textbf{相比现有方式节省或增加了哪些成本：}单次查询的全流程 API 调用成本约 \textbf{1.5 元}（含多模态逐页提取、质量评估与冲突归因等全部模型调用，按阿里云百炼按量计费），耗时 14--21 分钟全自动完成；对照人工整理同类查询约需 3.5 小时且仅能覆盖 34 条记录（见 5.1 对照组 1）。即以约两杯咖啡的显性成本换取一次“覆盖 30 个来源、百余条记录、逐条带溯源”的深度数据挖掘——折算到每条可用记录上的成本远低于人工逐篇摘录，且省去了人工整理最昂贵的部分：出处核对与单位换算。代价是端到端时延高于直接问答（21 分钟 vs 秒级），适合“研究者可等待、数据要可靠”的科研取数场景而非即时闲聊式提问。
\item \textbf{必须由研究者判断或人工审核的内容：}① \textbf{不同方法记录导致的不同取值}——同一物理量的多方法测量（如主序拟合法 vs Gaia 视差法）均真实有效时，保留哪条、全部纳入拟合、还是单独建成“多方法对照数据集”，由研究者按研究目的裁决；② \textbf{疑似错值的剔除}——绝大多数记录经自动质检拦截，但仍可能存在 1--2 条通过合理范围检查的错值（如成分混淆或口径混入），需要研究者结合领域知识在最终数据集中剔除或标注；③ 跨来源写法与单位取舍（无法自动互算的对数写法等）与系统送审的无法裁决项（见 2.4 承诺边界与 6.2 失败档案）。
\item \textbf{尚不能支持的数据来源、学科范围或复杂情况：}① \textbf{学科范围}：系统面向 SIMBAD 分类体系中 TAXONOMY OF STARS（恒星个体）、SETS OF STARS（恒星集合，星团/星协/双星等）、TAXONOMY OF GALAXIES（星系个体）、SETS OF GALAXIES（星系集合）四类天体提供完整性质体系；INTERSTELLAR MEDIUM（星际介质）、GRAVITATION（引力透镜等）、GENERAL SPECTRAL PROPERTIES（通用光谱属性）、BLENDS / ERRORS / NOT WELL DEFINED OBJECTS（混合体/错误/未明对象）四类不在当前性质库覆盖内，查询可能零产出或仅返回星表基础量；② \textbf{来源范围}：受学术出版版权生态制约，全文解析以 arXiv 预印本与开放获取渠道为主，封闭期刊正文暂无法直接获取（元数据仍可检索）；③ \textbf{复杂情况}：时变性质的连续监测序列、多星系统的成员级数值归属、“类名原型天体”的种群统计论文（论文主体讨论一整类天体而非目标本身）等仍需人工判断或尚不支持（见 6.2 失败档案），也是后续迭代的重点方向。
\end{itemize}

% ==================== 第 7 章 代码复现与提交材料 ====================
\section{代码复现与提交材料}

\subsection{最终交付内容}

\begin{longtable}{@{}p{4.2cm}p{10.2cm}@{}}
\toprule
交付项目 & 实际提交内容或入口 \\
\midrule
源代码、环境、依赖与运行方法 & \emph{（随提交材料附上：仓库结构、一键运行脚本、依赖与 README 入口）} \\
Qwen 模型、调用方式及凭证或截图 & \emph{（随提交材料附上：阿里云百炼 API Key 调用日志/控制台截图及计费凭证）} \\
代表案例的输入、来源、解析结果和修正记录 & \emph{（随提交材料附上：M45 案例数据目录与导出文件位置）} \\
结构化输出、字段说明和总体质量结果 & \emph{（随提交材料附上：最终导出的 JSON/CSV/Schema 说明路径）} \\
测试 API、示例请求与可交互前端 & \emph{（随提交材料附上：API 文档、示例请求、前端启动方式与样例库入口）} \\
可选演示视频 & \emph{（作为最佳：10 分钟以内演示视频链，随提交材料附上）} \\
\bottomrule
\end{longtable}

\subsection{提交前自检清单}

\begin{itemize}[leftmargin=1.2em,itemsep=1pt]
\item [$\square$] 技术方案文档 PDF 不超过 30 页，正文内容与本作品实际实现一致。
\item [$\square$] 代表案例真实呈现数据需求、来源查找、内容解析和数据统一过程；质量档案与版本迭代（V1/V2）如实呈现，未虚构版本。
\item [$\square$] 溯源清晰：数据来源、原始证据、模型推断和人工判断能够严格区分，结构化记录可以精准回溯到实际来源。
\item [$\square$] 对照客观：对照采用相同数据需求和评价口径，未省略错误、无法判断、未改善部分和适用边界。
\item [$\square$] 代码与服务可用：源代码、运行说明、Qwen 调用凭证、结构化成果、测试 API 和前端入口可以正常核验；受限数据、密钥和个人敏感信息已妥善脱敏处理。
\item [$\square$] 模板设问对应关系（P1--P20 覆盖、重复设问合并回答）见附录。
\end{itemize}
